% GENERATED FILE. Edit canonical lessons, then run npm run generate:books.
% canonical-source-hash: fnv1a64:477df6500c508b3c
% canonical-lessons: TA-C357-eninum, TA-C357-marupuram, TA-C357-tavira, TA-C357-atepol, TA-C357-verupatu

\chapter{On the Other Hand}
\label{ch:ta-on-the-other-hand}
\hlchaptermodality{357}

\begin{chapteropening}
\textbf{By the end of this chapter:} \emph{I can contrast two ideas and say how they differ.}

Complete the last lesson of chapter 357: I can contrast two ideas and say how they differ.
\end{chapteropening}

\section[\texorpdfstring{\ta{எனினும்} (eṉiṉum)}{எனினும் (eṉiṉum)}]{\ta{எனினும்} (eṉiṉum) — nevertheless; however}

\label{lesson:TA-C357-eninum}



\begin{quote}
Before the new one: say the Tamil for a compromise, then the Tamil for a heated argument.
\end{quote}

\subsection*{You'll want to know: \ta{எனினும்}}
\textbf{\ta{எனினும்}} — \emph{eṉiṉum} — \textquotedblleft{}nevertheless; however\textquotedblright{}.

\begin{quote}
\ta{மழை} \ta{பெய்தது}, \ta{எனினும்} \ta{போட்டி} \ta{நடந்தது}.
\end{quote}

\begin{quote}
\emph{maḻai peytatu, eṉiṉum pōṭṭi naṭantatu.}
\end{quote}

\begin{quote}
It rained. Nevertheless, the match took place.
\end{quote}

A formal written word; the everyday one is \ta{இருந்தாலும்}. It opens the second sentence, with a comma.

\begin{cousinweb}
\textbf{\ta{எனினும்}} (\emph{eṉiṉum}) is native: it is \emph{eṉ}, \textquotedblleft{}to say\textquotedblright{} (DEDR 868), with the concessive ending -\ta{இனும்}, \textquotedblleft{}even if\textquotedblright{}: \textquotedblleft{}even if it is said so\textquotedblright{}.
\end{cousinweb}

\begin{grammarlens}[title={What you've built}]
One more word for giving reasons and opinions.
\end{grammarlens}

\subsection*{Your turn}
\begin{itemize}
\raggedright
  \item \emph{Say it:} \emph{eṉiṉum}
  \item \emph{Say it:} \emph{eṉiṉum}, once more
  \item \emph{Say it:} \emph{maḻai peytatu, eṉiṉum pōṭṭi naṭantatu.} — It rained. Nevertheless, the match took place.
  \item \emph{Recall:} say the Tamil for a compromise, then the Tamil for a heated argument, then say \emph{eṉiṉum} again
\end{itemize}

\begin{tcolorbox}[breakable,colback=teal!4,colframe=teal!35!black,title={Before you move on}]
What does \ta{எனினும்} mean? (\textquotedblleft{}Nevertheless; however\textquotedblright{}.) Say it once more.
\end{tcolorbox}
\section[\texorpdfstring{\ta{மறுபுறம்} (maṟupuṟam)}{மறுபுறம் (maṟupuṟam)}]{\ta{மறுபுறம்} (maṟupuṟam) — on the other hand}

\label{lesson:TA-C357-marupuram}



\begin{quote}
Before the new one: say the Tamil for a heated argument, then the Tamil for nevertheless.
\end{quote}

\subsection*{You'll want to know: \ta{மறுபுறம்}}
\textbf{\ta{மறுபுறம்}} — \emph{maṟupuṟam} — \textquotedblleft{}on the other hand\textquotedblright{}.

\begin{quote}
\ta{இந்தக்} \ta{கைப்பேசி} \ta{மலிவானது}, \ta{மறுபுறம்} \ta{இது} \ta{மெதுவாக} \ta{வேலை} \ta{செய்கிறது}.
\end{quote}

\begin{quote}
\emph{intak kaippēci malivāṉatu, maṟupuṟam itu metuvāka vēlai ceykiṟatu.}
\end{quote}

\begin{quote}
This phone is cheap; on the other hand, it works slowly.
\end{quote}

\ta{ஒருபுறம்} … \ta{மறுபுறம்}, on the one hand … on the other.

\begin{cousinweb}
\textbf{\ta{மறுபுறம்}} (\emph{maṟupuṟam}) is native: \emph{maṟu}, \textquotedblleft{}another, other\textquotedblright{} (DEDR 4766), and \emph{puṟam}, \textquotedblleft{}side, outside\textquotedblright{} (DEDR 4333): the other side.
\end{cousinweb}

\begin{grammarlens}[title={What you've built}]
One more word for giving reasons and opinions.
\end{grammarlens}

\subsection*{Your turn}
\begin{itemize}
\raggedright
  \item \emph{Say it:} \emph{maṟupuṟam}
  \item \emph{Say it:} \emph{maṟupuṟam}, once more
  \item \emph{Say it:} \emph{intak kaippēci malivāṉatu, maṟupuṟam itu metuvāka vēlai ceykiṟatu.} — This phone is cheap; on the other hand, it works slowly.
  \item \emph{Recall:} say the Tamil for a heated argument, then the Tamil for nevertheless, then say \emph{maṟupuṟam} again
\end{itemize}

\begin{tcolorbox}[breakable,colback=teal!4,colframe=teal!35!black,title={Before you move on}]
What does \ta{மறுபுறம்} mean? (\textquotedblleft{}On the other hand\textquotedblright{}.) Say it once more.
\end{tcolorbox}
\section[\texorpdfstring{\ta{தவிர} (tavira)}{தவிர (tavira)}]{\ta{தவிர} (tavira) — except; apart from}

\label{lesson:TA-C357-tavira}



\begin{quote}
Before the new one: say the Tamil for nevertheless, then the Tamil for on the other hand.
\end{quote}

\subsection*{You'll want to know: \ta{தவிர}}
\textbf{\ta{தவிர}} — \emph{tavira} — \textquotedblleft{}except; apart from\textquotedblright{}.

\begin{quote}
\ta{ஞாயிறு} \ta{தவிர} \ta{எல்லா} \ta{நாட்களிலும்} \ta{கடை} \ta{திறந்திருக்கும்}.
\end{quote}

\begin{quote}
\emph{ñāyiṟu tavira ellā nāṭkaḷilum kaṭai tiṟantirukkum.}
\end{quote}

\begin{quote}
The shop is open every day except Sunday.
\end{quote}

It follows the noun: X \ta{தவிர}, except X; \ta{அதைத்} \ta{தவிர}, apart from that.

\begin{cousinweb}
\textbf{\ta{தவிர}} (\emph{tavira}) is native Dravidian: it is the infinitive of \emph{tavir}, \textquotedblleft{}to abstain, leave, avoid\textquotedblright{} (DEDR 3113), the verb \ta{தவிர்} taught already: \textquotedblleft{}leaving aside\textquotedblright{}.
\end{cousinweb}

\begin{grammarlens}[title={What you've built}]
One more word for giving reasons and opinions.
\end{grammarlens}

\subsection*{Your turn}
\begin{itemize}
\raggedright
  \item \emph{Say it:} \emph{tavira}
  \item \emph{Say it:} \emph{tavira}, once more
  \item \emph{Say it:} \emph{ñāyiṟu tavira ellā nāṭkaḷilum kaṭai tiṟantirukkum.} — The shop is open every day except Sunday.
  \item \emph{Recall:} say the Tamil for nevertheless, then the Tamil for on the other hand, then say \emph{tavira} again
\end{itemize}

\begin{tcolorbox}[breakable,colback=teal!4,colframe=teal!35!black,title={Before you move on}]
What does \ta{தவிர} mean? (\textquotedblleft{}Except; apart from\textquotedblright{}.) Say it once more.
\end{tcolorbox}
\section[\texorpdfstring{\ta{அதேபோல்} (atēpōl)}{அதேபோல் (atēpōl)}]{\ta{அதேபோல்} (atēpōl) — likewise; in the same way}

\label{lesson:TA-C357-atepol}



\begin{quote}
Before the new one: say the Tamil for on the other hand, then the Tamil for except.
\end{quote}

\subsection*{You'll want to know: \ta{அதேபோல்}}
\textbf{\ta{அதேபோல்}} — \emph{atēpōl} — \textquotedblleft{}likewise; in the same way\textquotedblright{}.

\begin{quote}
\ta{அவன்} \ta{கடினமாகப்} \ta{படித்தான்}, \ta{அதேபோல்} \ta{அவன்} \ta{தங்கையும்} \ta{படித்தாள்}.
\end{quote}

\begin{quote}
\emph{avaṉ kaṭiṉamākap paṭittāṉ, atēpōl avaṉ taṅkaiyum paṭittāḷ.}
\end{quote}

\begin{quote}
He studied hard, and likewise so did his younger sister.
\end{quote}

\ta{அதே}, that same, + \ta{போல்}, like: \ta{அதேபோல்} \ta{செய்}, do it the same way.

\begin{cousinweb}
\textbf{\ta{அதேபோல்}} (\emph{atēpōl}) is native: \emph{atē}, \textquotedblleft{}that very\textquotedblright{}, from \emph{atu}, \textquotedblleft{}that\textquotedblright{} (DEDR 1), and \emph{pōl}, \textquotedblleft{}to resemble; like\textquotedblright{} (DEDR 4597).
\end{cousinweb}

\begin{grammarlens}[title={What you've built}]
One more word for giving reasons and opinions.
\end{grammarlens}

\subsection*{Your turn}
\begin{itemize}
\raggedright
  \item \emph{Say it:} \emph{atēpōl}
  \item \emph{Say it:} \emph{atēpōl}, once more
  \item \emph{Say it:} \emph{avaṉ kaṭiṉamākap paṭittāṉ, atēpōl avaṉ taṅkaiyum paṭittāḷ.} — He studied hard, and likewise so did his younger sister.
  \item \emph{Recall:} say the Tamil for on the other hand, then the Tamil for except, then say \emph{atēpōl} again
\end{itemize}

\begin{tcolorbox}[breakable,colback=teal!4,colframe=teal!35!black,title={Before you move on}]
What does \ta{அதேபோல்} mean? (\textquotedblleft{}Likewise; in the same way\textquotedblright{}.) Say it once more.
\end{tcolorbox}
\section[\texorpdfstring{\ta{வேறுபடு} (vēṟupaṭu)}{வேறுபடு (vēṟupaṭu)}]{\ta{வேறுபடு} (vēṟupaṭu) — to differ; to vary}

\label{lesson:TA-C357-verupatu}



\begin{quote}
Before the new one: say the Tamil for except, then the Tamil for likewise.
\end{quote}

\subsection*{You'll want to know: \ta{வேறுபடு}}
\textbf{\ta{வேறுபடு}} — \emph{vēṟupaṭu} — \textquotedblleft{}to differ; to vary\textquotedblright{}.

\begin{quote}
\ta{இந்த} \ta{இரண்டு} \ta{கருத்துகளும்} \ta{பல} \ta{வழிகளில்} \ta{வேறுபடுகின்றன}.
\end{quote}

\begin{quote}
\emph{inta iraṇṭu karuttukaḷum pala vaḻikaḷil vēṟupaṭukiṉṟaṉa.}
\end{quote}

\begin{quote}
These two views differ in many ways.
\end{quote}

Past (I …): \ta{வேறுபட்டேன்}, \emph{vēṟupaṭṭēṉ}.

From what: X-\ta{இலிருந்து} \ta{வேறுபடு}. The noun is \ta{வேறுபாடு}, a difference.

\begin{cousinweb}
\textbf{\ta{வேறுபடு}} (\emph{vēṟupaṭu}) is native Dravidian: DEDR 5548 lists Tamil \emph{vēṟu}, \textquotedblleft{}other, different\textquotedblright{}, and \emph{vēṟu-paṭu}, \textquotedblleft{}to become different\textquotedblright{}, with \emph{paṭu}, \textquotedblleft{}to undergo\textquotedblright{}.
\end{cousinweb}

\begin{grammarlens}[title={What you've built}]
One more word for giving reasons and opinions.
\end{grammarlens}

\subsection*{Your turn}
\begin{itemize}
\raggedright
  \item \emph{Say it:} \emph{vēṟupaṭu}
  \item \emph{Say it:} \emph{vēṟupaṭu}, once more
  \item \emph{Say it:} \emph{inta iraṇṭu karuttukaḷum pala vaḻikaḷil vēṟupaṭukiṉṟaṉa.} — These two views differ in many ways.
  \item \emph{Recall:} say the Tamil for except, then the Tamil for likewise, then say \emph{vēṟupaṭu} again
\end{itemize}

\begin{tcolorbox}[breakable,colback=teal!4,colframe=teal!35!black,title={Before you move on}]
What does \ta{வேறுபடு} mean? (\textquotedblleft{}To differ; to vary\textquotedblright{}.) Say it once more.
\end{tcolorbox}
